\documentclass[11pt]{article}
\usepackage[a4paper,margin=27mm]{geometry}
\usepackage{amsmath,amssymb,amsthm,hyperref}
\newtheorem{theorem}{Theorem}
\newtheorem{corollary}{Corollary}
\title{A triple-last half-word gives four-wise fair dice}
\author{Research note: exact palindrome criterion and small constructions}
\date{30 September 2026}
\begin{document}
\maketitle

Say that a word is triple-last fair if on every triple of letters,
each die is largest with probability $1/3$.

\begin{theorem}\label{lift}
Let $V$ have $h$ occurrences of each of its letters. The palindrome
$W=V^{\rm rev}V$ is three-wise fair if and only if $V$ is triple-last
fair. Whenever these conditions hold, $W$ is also four-wise fair.
No pairwise fairness assumption on $V$ is needed.
\end{theorem}
\begin{proof}
Fix three letters $i,j,k$. For an occurrence of $i$ in $V$, let
$a,b$ be the numbers of earlier $j$'s and $k$'s. The corresponding
two occurrences in $W$ have prefix counts $(h-a,h-b)$ and
$(h+a,h+b)$. Hence their contributions to the count of selections
in which $i$ is largest sum to
\[
 (h-a)(h-b)+(h+a)(h+b)=2h^2+2ab.
\]
Summing over all $h$ occurrences of $i$ gives
\[
 L_i(W)=2h^3+2L_i(V).
\]
Therefore $L_i(W)=(2h)^3/3$ is equivalent to $L_i(V)=h^3/3$.
A palindrome gives the same count to an order and its reversal.
For a triple, the three equal last-place probabilities, together
with these reversal equalities, force all six order probabilities
to equal $1/6$. This proves the equivalence.

For completeness, the standard even-degree lifting argument applies
to any three-wise fair palindrome. In the completed free associative
algebra, give a word its signature $S=\prod\exp(X_{\rm letter})$.
Coefficients of words with distinct letters are exactly scattered
subword counts. The Baker--Campbell--Hausdorff logarithm $\log S$
is a Lie series. Reversal acts on a homogeneous Lie term of degree
$d$ by $(-1)^{d-1}$. Since the signature of a palindrome is fixed
by reversal, its even-degree logarithmic terms vanish. Three-wise
fairness makes its multilinear logarithmic terms of degrees two
and three vanish as well. Thus the coefficient of any four distinct
letters in $S$ comes solely from the fourth power of its degree-one
term, divided by $4!$. This is exactly four-wise fairness. Terms
with repeated letters cannot contribute to such a coefficient.
The same lifting principle is recorded in
\path{../partial_fairness/three_four_wise_near_linear.tex}.
\end{proof}

For the opposite orientation $VV^{\rm rev}$, replace ``triple-last''
by ``triple-first'' in Theorem~\ref{lift}.

\begin{corollary}
The minimum common face counts satisfy
\[
 M_4(5)=M_4(6)=12.
\]
\end{corollary}
\begin{proof}
The exact certificates
\path{../partial_fairness/triple_winner_half_n5_M6_seed0.json} and
\path{../partial_fairness/triple_winner_half_n6_M6_seed0.json}
give six-face triple-last fair half-words. In particular, the
five-letter half-word is
\[
 001002344321324114234231132400.
\]
The six-letter half-word is
\[
 010234432155042311324550055324114230.
\]
For every triple and every choice of its last label, each half-word
has exactly $6^3/3=72$ relevant selections. Their reverse-then-forward
palindromes therefore give twelve-face four-wise fair dice.
The lower bound follows from $M_4(n)\ge M_3(n)=12$ for $n=5,6$,
proved in \path{cloning_three_wise_dice.tex} using the six-face
extension obstruction.
\end{proof}

The independent verifier \path{verify_four_wise_palindromes.py}
enumerates index subsets literally, without reusing the constructor's
dynamic program. It checks every ordered pair, triple and quadruple:
their counts are respectively $72,288,864$. There are $120$ quadruple
orders in the five-letter case and $360$ in the six-letter case.
The verifier also checks the half-word last-place counts and the
exact orientation $W=V^{\rm rev}V$.
\end{document}
